\section{Excluding the residual allocation (2,2,1)}
\label{sec:exclude-two-two-one}

Continue with a finite nonempty strongly connected all-deficient graph
minimal under deletion of arbitrary nonempty arc sets.
Let $\mathcal B$ be its cross-witness graph: two vertices are adjacent
when they occur as opposite directional failure witnesses of a bad
missing pair. Every edge of $\mathcal B$ is an original missing pair.

\begin{lemma}[Large two-unit vertices have one witness partner]
\label{lem:eta-two-witness-degree}
If $\eta(x)=2$ and $t(x)>2$, then $d_{\mathcal B}(x)\le1$.
For its sole witness partner $c$, when one exists, $c\to x$ is a
convenient orientation.
\end{lemma}
\begin{proof}
If $P(x)=\varnothing$, the identity
$2=2\mu(x)+g(x)-1$ and $g(x)\in\{1,2\}$ force
$\mu(x)=1$, so there is only one missing partner.
That partner is conveniently directed into $x$ because every vertex
reaches $x$ within two steps.
Otherwise the local residual and unprotected-target bounds force
$s=0,j=1$ when $t(x)>2$. The exceptional missing-partner set is
$B_x=\{q\}$, and Lemma~\ref{lem:partner-head-budget} has zero
remaining head budget. Thus every missing partner $r\ne q$ satisfies
$N_D^+(r)\cap U_D(x)=\varnothing$.
A cross-witness partner necessarily has a nonempty such intersection,
so it must be $q$. That lemma also gives convenience of $q\to x$.
\end{proof}

\begin{lemma}[A two-unit leaf in a source star]
\label{lem:eta-two-leaf-completion}
Suppose the cross-witness graph is a star with center $c$, and $v$ is
a leaf with $\eta(v)=2$. Orient every bad missing pair in the direction
whose only failure source is $c$. Canonically orient good pairs using
a reference extension which places all good missing partners of $v$
before $v$. The resulting tournament $T$ keeps $v$ deficient.
If $t(v)>2$, the completion also adds the first arc $c\to v$.
\end{lemma}
\begin{proof}
Each bad pair has one directional failure set equal to $\{c\}$:
otherwise opposite noncentral witnesses would create an edge between
leaves of the star. This justifies the prescribed orientation.
Since $cv$ is missing, a selected bad direction cannot be an outgoing
arc from $v$: failure at $c$ would require the original arc $c\to v$.
Thus all added first neighbours of $v$ come from good pairs.

The reference extension exists, because every missing partner of $v$
is incomparable with $v$ in the protected order. If $t(v)\le2$, the
nonselected-source budget from
Lemma~\ref{lem:single-selected-source-completion} keeps $v$ deficient:
with no added first neighbour its gap is unchanged, and with at least
one it is at least $g_D(v)+2-t(v)\ge1$.

For $t(v)>2$ and $P(v)=\varnothing$, one has $\mu(v)=1$ and $g(v)=1$,
so the sole missing partner is $c$. The direction $c\to v$ is
convenient and is selected by the reference extension. There are no
added first neighbours at $v$, and its original first intermediates
cannot create a far head because the selected bad directions fail
only at $c$.

For $t(v)>2$ and $P(v)\ne\varnothing$, the preceding lemma identifies
the unique exceptional partner as $c$. Thus $s=0,j=1$ and the local
envelope saturates:
\[
 E=(\mathcal M(v)\setminus\{c\})\,\dot\cup\,\partial^2P(v)
   =N_D^{++}(v).
\]
Every vertex in $P(v)\cup\partial P(v)\cup\partial^2P(v)$ is
originally reachable from $v$ within two steps. The convenient
direction $c\to v$ is selected, so any added first neighbour
$r$ lies in $\mathcal M(v)\setminus\{c\}\subseteq E$.
Because all good partners were placed before $v$, an added good
direction $v\to r$ means that $r\to v$ is not convenient. Therefore
there is an original $p\in P(v)$ with $p\to r$.
In particular $p\ne c$, since $c$ is outside $P(v)$.
This argument applies to the actual added-first set; it does not
assume that every nonexceptional partner has residual zero or belongs
to $\partial P(v)$.

Let $y\in U_D(v)$. An original $r\to y$ is impossible by saturation.
Moreover $p$ cannot reach $y$ within two steps: all such paths from
$P(v)$ stay in the displayed reachable union.
If $ry$ is good, $r\to y$ is nonconvenient because it fails at $p$.
If $ry$ is bad, a selected direction $r\to y$ would likewise fail at
$p\ne c$, contradicting the choice of $c$ as its only failure source.
Thus the completion never directs $r$ toward $y$.
Original first intermediates of $v$ are also safe because bad
directions fail only at $c$.
Writing $K=N_T^+(v)\setminus N_D^+(v)$, we obtain
\[
 K\subseteq N_D^{++}(v),\qquad
 N_T^{++}(v)=N_D^{++}(v)\setminus K,
 \qquad g_T(v)=g_D(v)+2|K|>0.
\]
The convenience and selection of $c\to v$ also prove the final
assertion in both large-budget cases.
\end{proof}

\begin{theorem}\label{thm:no-two-two-one}
The positive residual allocation $(2,2,1)$ is impossible.
\end{theorem}
\begin{proof}
Let the positive vertices be $u,v,w$, with residuals $2,2,1$.
Every other vertex has $t=0$, while
Corollary~\ref{cor:eta-one-budget} gives $t(w)\le1$.
All directional failure sources lie in these three vertices.
The cross-witness graph $\mathcal B$ is nonempty.
Every missing pair inside this three-vertex set is good: a bad pair
would require two further distinct failure witnesses within the set.

If $\mathcal B$ is a triangle, the two residual-two vertices each have
two witness partners, and Lemma~\ref{lem:eta-two-witness-degree}
forces $t(u),t(v)\le2$. All source thresholds
$t(x)\le g(x)+1$ now hold, so the generic cross-witness transfer
Lemma~\ref{lem:cross-witness-source-transfer} gives a Seymour vertex.

Suppose $\mathcal B$ has one edge incident to the unit vertex $w$.
The actual directional failure sources are just the endpoints of
that edge; Theorem~\ref{thm:two-sources-one-unit} applies.
Thus the remaining one-edge possibility joins $u$ and $v$.
Orient all bad pairs with $u$ as their sole selected failure source,
and apply Lemma~\ref{lem:eta-two-leaf-completion} to the leaf $v$.
The generic single-selected-source completion budget keeps the
nonselected unit $w$ deficient, and keeps all zero-residual vertices
deficient. Only $u$ can be a Seymour vertex of this fixed tournament.

Next let $\mathcal B$ be a two-edge path centered at a residual-two
vertex, say $u$. Select $u$ as the sole failure source of every bad
direction, and protect the other residual-two leaf $v$ by
Lemma~\ref{lem:eta-two-leaf-completion}. The nonselected unit leaf
$w$ remains deficient by its one-unit budget. Again all vertices
except possibly $u$ remain deficient in the same tournament.

Finally let the path be centered at the unit vertex $w$. Select $w$
as the sole failure source. If both heavy leaves satisfy $t\le2$,
the generic nonselected-source budget keeps both deficient for any
canonical completion. Only $w$ can then be Seymour.
Otherwise choose a heavy leaf with $t>2$, say $v$, and use the
reference extension of Lemma~\ref{lem:eta-two-leaf-completion}.
This keeps $v$ deficient and adds the good first arc $w\to v$.
The center $w$ has no incident bad pair, since a selected direction
cannot fail at its own endpoint. Its new first neighbours are
therefore good ones. The center's general completion budget gives
\[
 g_T(w)\ge g_D(w)+2|K_w|-t(w)\ge g_D(w)+1>0,
\]
because $|K_w|\ge1$ and $t(w)\le1$.
Thus only the other heavy leaf $u$ can be Seymour.

These cases exhaust all nonempty graphs on three vertices.
In every non-transfer case we produced one fixed completed tournament
with at most one Seymour vertex. It has no sink, because every
original vertex had positive out-degree. Havet--Thomasse's theorem
provides at least two distinct Seymour vertices in such a tournament
\cite[Theorem~2]{HavetThomasse2000}, a contradiction.
\end{proof}
